\documentclass[10pt]{article}
\usepackage[T1]{fontenc}
\usepackage[utf8]{inputenc}
\usepackage[french]{babel}
\usepackage{newpxtext,newpxmath}
\usepackage[a4paper,top=1.5cm,bottom=1.35cm,left=1.55cm,right=1.55cm]{geometry}
\usepackage{microtype}
\usepackage{xcolor}
\usepackage{titlesec}
\usepackage[most]{tcolorbox}
\usepackage{enumitem}
\usepackage{array}
\usepackage{booktabs}
\usepackage{fancyhdr}

\definecolor{navy}{HTML}{1F3A5F}
\definecolor{grey}{HTML}{54606E}
\definecolor{green}{HTML}{2C6E49}
\definecolor{amber}{HTML}{9C6A16}
\definecolor{red}{HTML}{A6322E}
\definecolor{plum}{HTML}{7A4EA8}
\definecolor{boxgray}{HTML}{F3F5F8}
\definecolor{boxred}{HTML}{FCF2F1}
\definecolor{boxgreen}{HTML}{F1F7F3}
\definecolor{boxplum}{HTML}{F7F3FB}
\definecolor{boxedge}{HTML}{D9E0EA}

\setlength{\parindent}{0pt}
\setlength{\parskip}{2.6pt}
\linespread{1.0}

\pagestyle{fancy}\fancyhf{}
\lhead{\small\color{grey}Scella}
\rhead{\small\color{grey}Couples : éléments, classement, usages -- 11 août 2026}
\cfoot{\small\color{grey}\thepage}
\renewcommand{\headrulewidth}{0.3pt}

\titleformat{\section}{\color{navy}\large\bfseries}{}{0pt}{}
\titlespacing{\section}{0pt}{9pt}{3pt}
\titleformat{\subsection}{\color{navy}\bfseries}{}{0pt}{}
\titlespacing{\subsection}{0pt}{7pt}{1.5pt}

\newtcolorbox{these}{colback=boxgray,colframe=boxedge,boxrule=0.5pt,arc=2.5pt,
  left=8pt,right=8pt,top=4.5pt,bottom=4.5pt}
\newtcolorbox{alerte}{colback=boxred,colframe=red!55,boxrule=0.8pt,arc=2.5pt,
  left=8pt,right=8pt,top=4.5pt,bottom=4.5pt}
\newtcolorbox{acquis}{colback=boxgreen,colframe=green!45,boxrule=0.8pt,arc=2.5pt,
  left=8pt,right=8pt,top=4.5pt,bottom=4.5pt}
\newtcolorbox{piste}{colback=boxplum,colframe=plum!45,boxrule=0.8pt,arc=2.5pt,
  left=8pt,right=8pt,top=4.5pt,bottom=4.5pt}

\newcommand{\best}[1]{{\color{green}\textbf{#1}}}
\newcommand{\pire}[1]{{\color{red}\textbf{#1}}}
\newcommand{\sz}[1]{{\small\color{grey}\textsf{#1}}}
\newcommand{\tz}[1]{{\footnotesize\color{grey}\textsf{#1}}}
\newcommand{\code}[1]{\texttt{\small #1}}
\newcommand{\x}{$\times$}

\begin{document}

\begin{center}
{\Huge\bfseries\color{navy} Scella}\\[4pt]
{\large\color{grey}\itshape Les couples : de quoi ils sont faits, ce qu'ils valent, contre qui ils tiennent}\\[4pt]
{\small\color{grey} 11 août 2026 \; · \; anatomie, classement intrinsèque et modèle de menace}
\end{center}
\vspace{2pt}

\begin{these}
\textbf{Ce que ce document ajoute aux deux précédents.} Le premier disait \emph{quels} couples existent,
le second \emph{dans quel ordre} ils valent. Celui-ci les ouvre : chaque couple est décomposé en ses
deux éléments \textbf{plus son liant} -- la pièce qui fait la multiplication, et celle qu'on oublie
toujours en décrivant une architecture. Puis chacun est confronté à la seule question qui donne un
périmètre plutôt qu'une promesse : \textbf{contre qui tient-il, et contre qui ne tient-il pas ?} Le
résultat d'ensemble est que ce portefeuille sert à \textbf{rendre un processus contrôlable}, et pas
encore à le rendre contrôlable \emph{par un adversaire complet}. Deux fils séparent les deux états,
et ils sont nommés.
\end{these}

\section{1 · Pourquoi le liant, et pas seulement les deux moitiés}

Un couple n'est pas une paire de composants voisins : c'est deux éléments \textbf{plus le mécanisme
qui les fait produire ensemble une propriété qu'aucun n'a seul}. Ce troisième terme -- le
\emph{liant} -- est presque toujours absent des descriptions d'architecture, parce qu'un diagramme
montre des boîtes et des flèches, pas des ordres d'exécution ni des contraintes d'appariement.

C'est pourtant lui qui porte toute la valeur, et le tableau du §2 le montre bloc par bloc. Le
caviardage PII et le scellement du reçu existent tous deux dans quantité de systèmes ; ce qui est rare
est que le premier s'exécute \emph{avant} le second. Une topologie déclarée et des enregistrements
d'appels sont ordinaires ; ce qui mord est que leur appariement soit \emph{injectif}. Une chaîne HMAC
et une signature Ed25519 sont deux primitives banales ; ce qui compte est qu'elles soient
\emph{orthogonales}, au point qu'une forgerie exige les deux clés.

Retenir cette lecture change la manière d'auditer : quand une propriété composée disparaît, ce n'est
presque jamais parce qu'un des deux éléments a été retiré -- ils sont visibles, quelqu'un le
remarquerait. C'est parce que le liant a été affaibli en silence : un appel déplacé, une contrainte
relâchée, une vérification rendue optionnelle.

\section{2 · Les éléments, couple par couple}

\renewcommand{\arraystretch}{1.25}
\begin{center}\footnotesize
\begin{tabular}{@{}>{\raggedright\arraybackslash}p{2.75cm} >{\raggedright\arraybackslash}p{4.05cm} >{\raggedright\arraybackslash}p{4.05cm} >{\raggedright\arraybackslash}p{5.15cm}@{}}
\toprule
\textbf{Le couple} & \textbf{Élément A} & \textbf{Élément B} & \textbf{Le liant -- ce qui fait la propriété} \\
\midrule
\textbf{1 · Artefact \x{} re-prouveur} &
  L'artefact scellé : \code{meta\_steps.jsonl} + \code{manifest.json} &
  Le re-prouveur : SHA-256, HMAC et parseur JSON \textbf{écrits à la main} &
  \textbf{L'absence de chaîne d'approvisionnement.} \code{cargo tree -e normal} montre le binaire
  \textbf{seul} : rien à qui faire confiance entre l'auditeur et l'octet \\
\textbf{2 · Topologie \x{} reçu} &
  La topologie \textbf{déclarée} -- la liste des appels annoncés &
  Un enregistrement scellé \textbf{par appel réellement effectué} &
  \textbf{L'appariement injectif.} Un n\oe ud déclaré $\leftrightarrow$ un record, jamais deux vers le
  même. Sans lui, 3 déclarés + 1 record $\Rightarrow$ « Consistent » \\
\textbf{3 · Reçu v4 \x{} ancre} &
  \code{participant\_ids} \textbf{ordonnés}, \code{declared\_config\_hashes} alignés, opérateur,
  $\alpha$, $\gamma$ &
  Le journal append-only : racine \textbf{et} nombre de pas &
  \textbf{L'antériorité de la publication.} Le préfixe d'une chaîne HMAC valide est valide ; seule une
  racine publiée \emph{avant} expose la coupure -- sans clé \\
\textbf{4 · Canari \x{} participants} &
  L'estimateur de Lipschitz empirique, boîte noire, bisection pour les discontinuités &
  Le participant efe-core réel et sa porte de démarrage &
  \textbf{Le fail-closed.} La mesure ne produit pas un rapport : elle produit un \emph{refus de
  servir} \\
\textbf{5= · PII \x{} scellement} &
  Le caviardage selon la policy tenant -- détection \textbf{locale}, sans réseau &
  Le hash du prompt scellé dans la chaîne &
  \textbf{L'ordre.} Caviarder \emph{avant} tout usage du texte $\Rightarrow$ le sceau couvre ce qui est
  réellement sorti \\
\textbf{5= · Reçu \x{} coût} &
  Le plafonnement : cap tenant \textbf{puis} plafond serveur, inconditionnel &
  Le coût par appel, scellé dans chaque enregistrement &
  \textbf{La borne en amont du sceau.} Le reçu enregistre une dépense que quelque chose avait
  \emph{déjà} bornée \\
\textbf{7 · R3 \x{} HMAC} &
  La chaîne HMAC tenant, qui scelle le record \textbf{co-signature incluse} &
  La signature Ed25519 du n\oe ud sur \code{prompt\_hash} $\Vert$ \code{output\_hash} &
  \textbf{L'orthogonalité des deux clés} : forger exige \textbf{les deux}. Plus la non-élévation :
  signature invérifiable $\Rightarrow$ redescente automatique \\
\textbf{8 · Revendication \x{} chute} &
  Les hops de reformulation, chacun portant son moteur (SLM / règle / dégradé) &
  Le prédicat « tous les hops sont SLM », calculé \textbf{depuis la sortie scellée} &
  \textbf{Le calcul depuis la source scellée.} La vue ne peut pas diverger de l'attesté ; une seule
  dégradation fait tomber la revendication \\
\bottomrule
\end{tabular}
\end{center}

\section{3 · Le classement (rappel)}

Établi le 11 août sur les quatre lentilles intrinsèques -- maturité, rareté, adversaire,
démontrabilité -- après le retrait de la lentille d'objectif extérieur.

\renewcommand{\arraystretch}{1.1}
\begin{center}\small
\begin{tabular}{@{}r >{\raggedright\arraybackslash}p{6.3cm} c c >{\raggedright\arraybackslash}p{6.0cm}@{}}
\toprule
& \textbf{Le couple} & \textbf{Moy} & \sz{Éc.} & \textbf{État} \\
\midrule
\textbf{1} & Artefact scellé \x{} re-prouveur zéro-dépendance & \textbf{2,00} & 2 & servi, sans caveat relevé \\
\textbf{2} & Topologie déclarée \x{} reçu apparié injectivement & \textbf{2,50} & 2 & servi \\
\textbf{3} & Reçu v4 \x{} ancre & \textbf{3,00} & 4 & servi ; ancre détenue par l'opérateur \\
\textbf{4} & Canari H1 \x{} participants réels & \textbf{4,00} & 4 & servi ; échappatoire déclarée \\
5\,= & Caviardage PII \x{} scellement & 5,75 & 3 & servi ; sans policy, sortie en clair \\
5\,= & Reçu \x{} coût réel par appel & 5,75 & 3 & servi \\
7 & Co-signature R3 \x{} chaîne HMAC & 6,25 & \pire{7} & \textbf{vérifieur servi, producteur absent} \\
8 & Revendication \x{} sa propre chute & 6,75 & 1 & servi ; voisin d'une fissure \\
\bottomrule
\end{tabular}
\end{center}

\sz{Éc. = amplitude entre le meilleur et le pire rang obtenus par un même couple. Un écart élevé
signale l'instabilité, pas la médiocrité : le rang 7 est premier en rareté et dernier partout ailleurs.}

\section{4 · À quoi ça sert -- et surtout, contre qui ça tient}

C'est la formulation qui rend un usage discutable. « À quoi ça sert » sans modèle de menace produit une
promesse ; la même question avec un adversaire nommé produit un \textbf{périmètre}. La colonne de
droite est donc la plus importante des trois, et c'est celle qu'il faut lire en premier.

\renewcommand{\arraystretch}{1.25}
\begin{center}\footnotesize
\begin{tabular}{@{}r >{\raggedright\arraybackslash}p{4.55cm} >{\raggedright\arraybackslash}p{5.35cm} >{\raggedright\arraybackslash}p{5.35cm}@{}}
\toprule
& \textbf{Le problème résolu} & \textbf{Tient contre} & \textbf{Ne tient PAS contre} \\
\midrule
\textbf{1} & Faire contrôler une affirmation par quelqu'un qui ne vous accorde rien -- ni à vous, ni à
  vos dépendances &
  un opérateur qui ment sur son propre outillage ; un auditeur \textbf{sans réseau} ; une vérification
  dans dix ans &
  un détenteur de clé qui fabrique la chaîne \textbf{entière} dès l'origine \\
\textbf{2} & Prouver qu'une consultation multi-sources a eu lieu \textbf{comme annoncée} &
  un opérateur qui prétend un panel et ne fait qu'un appel ; le même fournisseur compté trois fois &
  un manifeste qui déclare \emph{honnêtement} un panel étroit -- il prouve la conformité à la
  déclaration, pas son ambition \\
\textbf{3} & Empêcher qu'on ne publie que la moitié qui arrange &
  la troncature rétroactive, \textbf{détectable sans aucune clé} &
  un opérateur qui n'ancre pas, ou qui réécrit son journal avant publication -- l'ancre est chez lui
  \sz{(c'est 3b)} \\
\textbf{4} & Refuser de servir quand l'hypothèse de sûreté du système est falsifiée &
  une régression silencieuse de l'opérateur ; une borne théorique optimiste -- elle \textbf{a été
  réfutée}, $\hat{L}=0{,}69$ contre $\beta_1 \approx 0{,}17$ annoncé &
  un manifeste qui désactive le canari : l'échappatoire est déclarée et bruyante, mais elle existe \\
\textbf{5=} & Répondre « qu'avez-vous envoyé à ce tiers ? » par une \textbf{preuve} et non par une
  affirmation &
  une fuite par oubli ; la disparition \emph{silencieuse} du caviardage quand la capacité n'est pas
  compilée &
  un tenant \textbf{sans policy} : le prompt sort en clair, et c'est documenté \\
\textbf{5=} & Une trace de dépense opposable, bornée en amont &
  le déni de service par le coût ; un client qui dimensionne la dépense d'autrui ; un dépassement de
  conversion &
  un opérateur qui trafique sa propre table de prix \\
\textbf{7} & Savoir \textbf{qui} a produit une sortie, prouvable sans détenir le secret du producteur &
  \emph{le concevrait contre} un n\oe ud de calcul distant qui nie avoir produit X &
  \pire{rien, aujourd'hui} -- aucun producteur, registre vide, aucun enregistrement n'est R3 \\
\textbf{8} & Empêcher une revendication de survivre à sa propre dégradation &
  un relais qui retombe silencieusement sur des règles au lieu du petit modèle &
  la souveraineté du relais : l'URL du petit modèle n'est pas contrainte au local \sz{(faux couple F1)} \\
\bottomrule
\end{tabular}
\end{center}

\section{5 · Les trois familles -- ce que l'ensemble fait réellement}

Les huit couples ne forment pas une liste plate. Ils remplissent trois fonctions distinctes, et cette
partition en dit plus long que le classement lui-même.

\textbf{Famille A -- prouver ce qui s'est passé à quelqu'un qui ne vous croit pas.} \sz{(couples 1, 2, 3)}
C'est le c\oe ur, et c'est aussi le haut du classement -- ce n'est pas une coïncidence, c'est le
résultat. Trois des quatre premiers font la même chose sous trois angles : rendre une affirmation
contrôlable \emph{malgré} l'intérêt de celui qui l'émet. Aucun ne demande qu'on croie l'opérateur ; ils
demandent seulement qu'il ait laissé une trace avant de savoir qu'on la lui réclamerait.

\textbf{Famille B -- empêcher le système de mentir sur lui-même.} \sz{(couples 4, 8, et la règle de non-élévation du 7)}
Fonction différente, et plus rare. Il ne s'agit plus de convaincre un tiers mais d'installer des
mécanismes qui \textbf{échouent bruyamment} quand une hypothèse interne cesse d'être vraie : le canari
qui refuse de démarrer, la revendication qui tombe à la première dégradation, le barreau qui redescend
sans qu'on le lui demande. C'est la famille la plus caractéristique du dépôt, et la moins visible de
l'extérieur -- elle ne produit rien tant que tout va bien.

\textbf{Famille C -- borner et tracer ce qui sort et ce que ça coûte.} \sz{(les deux couples 5=)}
La plus ordinaire, et le classement le dit sans détour : ce sont les deux derniers du groupe servi.
Nécessaires, solides, non distinctifs. Les compter comme des atouts serait la première erreur de
présentation.

\section{6 · Ce que ce portefeuille sert, et la limite qui le borne}

La forme de problème que ces couples adressent est précise, et elle mérite d'être énoncée telle
quelle : \textbf{un processus automatisé, exécuté par une partie qui a un intérêt dans son résultat, et
dont un tiers doit pouvoir contrôler le déroulement après coup}. Toute situation de cette forme est un
usage possible -- audit d'un système de décision, expertise après litige, archivage réglementaire,
calcul délégué à un fournisseur, reproductibilité d'un protocole. Ce n'est pas une liste de marchés :
c'est la description d'une \emph{forme}, et elle est volontairement laissée à ce niveau.

\begin{alerte}
\textbf{La limite qui borne l'ensemble, et qu'il vaut mieux dire avant qu'après.} Le reçu est
\emph{tamper-evident pour un détenteur de clé}, pas \emph{trustless pour un tiers}. Il ne rend pas le
mensonge impossible : il en augmente le coût et en laisse la trace. Deux exceptions -- et ce sont
précisément les deux mieux classées : le \textbf{re-prouveur zéro-dépendance}, qui ne demande de
confiance à personne, et l'\textbf{ancre}, qui expose la troncature sans clé.
\end{alerte}

Autrement dit, et c'est la phrase à retenir : ce portefeuille sert aujourd'hui à \textbf{rendre un
processus contrôlable}, et pas encore à le rendre contrôlable \textbf{par un adversaire complet}. La
distance entre les deux états ne tient pas à une découverte à faire, mais à \textbf{deux fils
identifiés} -- un producteur de co-signature sur le chemin servi, et un ancrage tiers. Les deux ont
leurs matériaux déjà écrits et testés dans le dépôt.

\section{7 · L'élément porteur -- ce que la décomposition révèle}

Décomposer huit couples en seize éléments fait apparaître une chose qu'aucune des deux analyses
précédentes ne pouvait montrer : \textbf{ces éléments ne sont pas seize}. Plusieurs sont le même,
réutilisé -- et l'un d'eux porte presque tout.

\renewcommand{\arraystretch}{1.15}
\begin{center}\small
\begin{tabular}{@{}>{\raggedright\arraybackslash}p{5.1cm} c >{\raggedright\arraybackslash}p{9.6cm}@{}}
\toprule
\textbf{L'élément} & \textbf{Couples} & \textbf{Ce que sa perte emporterait} \\
\midrule
\textbf{La chaîne HMAC scellée} & \pire{7 / 8} &
  1, 2, 3, les deux 5=, 7 et 8 \textbf{simultanément}. C'est le point porteur unique du portefeuille \\
La policy tenant & 2 & les deux couples 5= -- caviardage \emph{et} plafond de coût en dépendent \\
La sortie scellée comme source unique & 2 & le couple 8, et l'appariement du couple 2 \\
La signature Ed25519 & 1 \sz{(+1 latent)} & le couple 7, et le chantier de signature de racine \\
\midrule
\textbf{L'estimateur de Lipschitz} & \best{1} &
  le couple 4 \textbf{seul} -- c'est le seul qui ne touche pas au reçu \\
\bottomrule
\end{tabular}
\end{center}

Le constat se vérifie en une ligne : \code{participant/local.rs}, qui porte le canari, ne contient
\textbf{aucune} occurrence de HMAC. Le couple 4 est le seul du portefeuille qui survivrait
intégralement à une faille dans le scellement -- parce qu'il ne prouve rien à personne : il refuse de
démarrer.

\begin{alerte}
\textbf{La conséquence, et elle réordonne les priorités.} Sept couples sur huit reposent sur le
\emph{même} élément. Une faille dans le scellement HMAC ne dégraderait pas le portefeuille, elle
l'\textbf{effondrerait d'un coup} -- y compris les couples les mieux classés, y compris ceux dont la
propriété semble indépendante. Cette concentration n'est pas un défaut de conception : c'est le prix
d'une architecture cohérente, où tout se ramène à un sceau unique. Mais elle doit être \emph{sue}, car
elle change ce qu'il faut protéger en priorité.
\end{alerte}

Et elle donne au rang 1 une justification qu'il n'avait pas encore. Si un seul élément porte sept
couples, alors la seule chose qui compte vraiment est de pouvoir \textbf{contrôler cet élément sans
dépendre de celui qui l'a produit}. C'est exactement ce que fait le re-prouveur zéro-dépendance -- et
c'est pourquoi il finit premier d'un classement où il n'était même pas engagé. Ce n'est pas un
amplificateur : c'est le \textbf{seul contrôle indépendant du point porteur unique}.

Ce qui rend d'autant plus gênant le fait, établi dans le document compagnon, qu'\textbf{aucun test ni
job d'intégration continue} ne garde la contrainte de dépendances nulles.

\begin{these}
\small\textbf{Portée et statut.} Les éléments, les liants et les emplacements sont des faits établis
contre le code du 11 août ; les colonnes « tient / ne tient pas » en découlent directement. Les
\textbf{notes} du classement, elles, sont un jugement structuré et discutable, pas une mesure --
distinction développée dans le document compagnon. Aucune cible d'usage n'est déclarée et ce document
n'en propose aucune : il décrit une forme de problème, pas un destinataire. La couche nano et
\code{efe-gold-test-gate} n'ont été examinés dans aucun des trois documents.
\end{these}

\end{document}
